%% Đề 7 -- de-giua-hoc-ki-1-toan-9-nam-2023-2024-phong-gddt-bac-son-lang-son.pdf
\deheader{ĐỀ SỐ 7}%
  {Đề kiểm tra Giữa Học kỳ — Năm học 2023--2024}%
  {Phòng GD\&ĐT | Chỉ phần Đại số}
\addcontentsline{toc}{section}{Đề 7 — Đề kiểm tra Giữa HK — Năm học 2023--2024}

\begin{de_cau}{Câu 1}{2,0 điểm}
\begin{enumerate}[label=\alph*)]
\item Tính: $\sqrt{75} - 2\sqrt{48} + \sqrt{27}$.
\item Rút gọn $P = \dfrac{x\sqrt{x}-1}{x-\sqrt{x}} - \dfrac{x\sqrt{x}+1}{x+\sqrt{x}} + \dfrac{x+1}{\sqrt{x}}$
với $x>0, x\neq 1$.
\end{enumerate}
\end{de_cau}

\begin{loi_giai_cau}
\textbf{a)} $5\sqrt{3} - 8\sqrt{3} + 3\sqrt{3} = 0$.

\textbf{b)}
\[P = \frac{(\sqrt{x}-1)(x+\sqrt{x}+1)}{\sqrt{x}(\sqrt{x}-1)} - \frac{(\sqrt{x}+1)(x-\sqrt{x}+1)}{\sqrt{x}(\sqrt{x}+1)} + \frac{x+1}{\sqrt{x}}\]
\[= \frac{x+\sqrt{x}+1}{\sqrt{x}} - \frac{x-\sqrt{x}+1}{\sqrt{x}} + \frac{x+1}{\sqrt{x}} = \frac{2\sqrt{x}+x+1}{\sqrt{x}} = \sqrt{x}+\frac{(\sqrt{x}+1)^2-2\sqrt{x}}{\sqrt{x}}.\]
Rút gọn: $P = \dfrac{x+2\sqrt{x}+1}{\sqrt{x}} = \dfrac{(\sqrt{x}+1)^2}{\sqrt{x}}$.
\end{loi_giai_cau}

\begin{de_cau}{Câu 2}{2,0 điểm}
\begin{enumerate}[label=\alph*)]
\item Giải phương trình $\sqrt{2x-1} = 3$.
\item Giải hệ phương trình 
\[ \begin{cases}5x+3y=11\\2x-y=4\end{cases} \].
\end{enumerate}
\end{de_cau}

\begin{loi_giai_cau}
\textbf{a)} ĐKXĐ: $x\geq\frac{1}{2}$.
$\sqrt{2x-1}=3\Rightarrow 2x-1=9\Rightarrow x=5$ (thỏa mãn).

\textbf{b)} Từ (2): $y=2x-4$. Thế vào (1): $5x+3(2x-4)=11\Rightarrow 11x=23\Rightarrow x=\frac{23}{11}$.
$y=\frac{46}{11}-4=-\frac{2}{11}$.

Nghiệm $(x;y)=\left(\dfrac{23}{11};-\dfrac{2}{11}\right)$.
\end{loi_giai_cau}

\begin{de_cau}{Câu 3}{2,0 điểm}
Tổng hai số bằng 28. Biết rằng $\frac{3}{4}$ số thứ nhất bằng $\frac{2}{3}$ số thứ hai.
Tìm hai số đó.
\end{de_cau}

\begin{loi_giai_cau}
Gọi hai số là $x$ và $y$, $x+y=28$ và $\dfrac{3x}{4}=\dfrac{2y}{3}$.

Từ điều kiện thứ hai: $9x=8y\Rightarrow y=\dfrac{9x}{8}$.

Thế vào điều kiện thứ nhất: $x+\dfrac{9x}{8}=28\Rightarrow \dfrac{17x}{8}=28\Rightarrow x=\dfrac{224}{17}\approx 13{,}18$.

$y=28-\dfrac{224}{17}=\dfrac{252}{17}\approx 14{,}82$.
\end{loi_giai_cau}

\begin{de_cau}{Câu 4}{1,0 điểm}
Cho $A = \dfrac{\sqrt{x}+1}{\sqrt{x}-1}$ với $x\geq 0, x\neq 1$. Tìm $x$ nguyên để $A$ nguyên.
\end{de_cau}

\begin{loi_giai_cau}
$A = \dfrac{\sqrt{x}+1}{\sqrt{x}-1} = 1 + \dfrac{2}{\sqrt{x}-1}$.

Để $A$ nguyên, $\dfrac{2}{\sqrt{x}-1}$ phải nguyên, tức $(\sqrt{x}-1)\in\{1,-1,2,-2\}$.

- $\sqrt{x}-1=1\Rightarrow x=4$ (thỏa, nguyên)
- $\sqrt{x}-1=2\Rightarrow x=9$ (thỏa, nguyên)
- $\sqrt{x}-1=-1\Rightarrow x=0$ (thỏa, nguyên)
- $\sqrt{x}-1=-2\Rightarrow \sqrt{x}=-1$ (loại)

Vậy $x\in\{0,4,9\}$.
\end{loi_giai_cau}
